\section{Full-energy capacity for every Lipschitz activation}
\label{sec:all-lipschitz-full-energy}
A scalar covering argument removes the bounded-variation condition on the
activation derivative, at the expense of additional logarithms in sample
size. The activation need only be globally Lipschitz.

For a set $F\subset[-1,1]^n$, write
$\mathfrak R(F)=n^{-1}\mathbb E_\epsilon\sup_{f\in F}\sum_i\epsilon_i f_i$
and let $d_n(f,g)=\|f-g\|_2/\sqrt n$. Let $\mathcal L$ be all scalar
one-Lipschitz functions on $[-1,1]$ which vanish at zero.

\begin{lemma}[An unknown scalar Lipschitz link]
\label{lem:unknown-link}
There is an absolute $C$ such that, for every symmetric nonempty
$F\subset[-1,1]^n$ and $n\ge2$,
\[
 \mathfrak R(\{\psi\circ f:f\in F,\ \psi\in\mathcal L\})
 \le C\log^{3/2}(2n)\mathfrak R(F)+C/\sqrt n.
\]
\end{lemma}
\begin{proof}
We may take closures in the compact cube. Let
$G(F)=n^{-1}\mathbb E_g\sup_{f\in F}\sum_i g_i f_i$, with independent
standard Gaussian coordinates. Conditioning on their magnitudes and
using the Rademacher contraction principle gives
\[
 G(F)\le\mathbb E\max_i|g_i|\,\mathfrak R(F)
       \le C\sqrt{\log(2n)}\,\mathfrak R(F).
\]
Gaussian Sudakov minoration, applied to a maximal separated subset,
therefore gives for $0<a\le2$
\[
 \log\mathcal N(a,F,d_n)
 \le CnG(F)^2/a^2
 \le Cn\log(2n)\mathfrak R(F)^2/a^2.
\]
The scalar class $\mathcal L$ has
$\log\mathcal N(a,\mathcal L,\|\cdot\|_\infty)\le C/a$ for
$0<a\le1$. One can discretize a grid of spacing proportional to $a$,
quantize function values at the same scale, and count only sequences
whose successive quantized increments have bounded integer magnitude;
the initial value is fixed at zero. Representatives may be chosen from
the original class in each nonempty cell.

Composition of these two covers yields
\[
 \log\mathcal N(3a,\mathcal L\circ F,d_n)
 \le Cn\log(2n)\mathfrak R(F)^2/a^2+C/a.
\]
The truncated Dudley entropy bound for Rademacher processes now implies
\[
 \mathfrak R(\mathcal L\circ F)
 \le4\alpha+\frac C{\sqrt n}\int_\alpha^2
          \sqrt{\log\mathcal N(a,\mathcal L\circ F,d_n)}\,da.
\]
Choose $\alpha=n^{-1/2}$. The first entropy term integrates as
$C\sqrt{\log(2n)}\mathfrak R(F)\log(2/\alpha)$, while the second
integrates to $C/\sqrt n$. This proves the claim. Gaussian Sudakov
minoration and the truncated Dudley bound are standard Gaussian and
Rademacher process inequalities~\cite{talagrand}; their metrics and normalizations have
been stated explicitly here.
\end{proof}

Let $\mathcal A_0$ be the closed absolute convex hull of the constant
function one and the $d$ coordinate functions on the input cube.
Recursively define
\[
 \mathcal A_\ell=\overline{\operatorname{aconv}}
   \bigl(\{1\}\cup\{\psi\circ f:f\in\mathcal A_{\ell-1},
                                     \psi\in\mathcal L\}\bigr).
\]
These are symmetric classes with envelope one. Absolute convexification
preserves the Rademacher supremum of a symmetric class. Adding the
constant atom increases its expected supremum by at most $1/\sqrt n$.
Lemma~\ref{lem:unknown-link} and the finite-coordinate bound consequently
show, on every input sample,
\begin{equation}
 \widehat{\mathfrak R}_n(\mathcal A_H)
 \le \frac{C^{H+1}\log^{3H/2}(2n)
                      [1+\sqrt{\log(2d+2)}]}{\sqrt n}.
 \label{eq:all-lipschitz-atom-rad}
\end{equation}
The same bound holds for worst-case empirical complexity.

Now let $\phi$ be globally $A$-Lipschitz, $A\ge1$, and put
$\beta=|\phi(0)|$. At a nonzero preactivation budget $t$, the normalized
profile
\[
 \psi_t(u)=\frac{\phi(tu)-\phi(0)}{At},\qquad -1\le u\le1,
\]
belongs to $\mathcal L$. Assign neuron budgets
\[
 r_j^1=\beta+A\sum_i|W^1_{ji}|,\qquad
 r_j^\ell=\beta+A\sum_i|W^\ell_{ji}|r_i^{\ell-1},\qquad
 P=\sum_j|v_j|r_j^H.
\]
A zero budget gives the zero feature. Otherwise its feature divided by
its budget is a convex combination of the appropriately signed constant atom and a
normalized profile applied to an absolute convex combination from the
preceding class. Induction proves $z_w\in P\mathcal A_H$.

For the usual weighted energy
$Q_N=(d/N)\|W^1\|_F^2+\sum_{\ell\ge2}\|W^\ell\|_F^2+N\|v\|^2$,
Cauchy--Schwarz and $\||W|\|_{\rm op}\le\|W\|_F$ give
\[
 \|r^1\|_2/\sqrt N\le\beta+\sqrt2A\sqrt{Q_N},\qquad
 \|r^\ell\|_2/\sqrt N
    \le\beta+A\sqrt{Q_N}\|r^{\ell-1}\|_2/\sqrt N.
\]
Hence, with $T=2+\beta+2A$ and $L=H+1$,
\[
 P\le T^H(1+Q_N)^{L/2}.
\]
Combining this with \eqref{eq:all-lipschitz-atom-rad} gives
\begin{equation}
 \widehat{\mathfrak R}_n(\{z_w:Q_N(w)\le q\})
 \le\frac{C_{H,A,\beta}\log^{3H/2}(2n)
          [1+\sqrt{\log(2d+2)}]}{\sqrt n}(1+q)^{L/2},
 \qquad |z_w|\le T^H(1+q)^{L/2}.
 \label{eq:all-lipschitz-full-energy-rad}
\end{equation}
The classes include all learner widths. For statistical separability,
use rational-parameter approximations of the actual networks, at each
integer width. Continuity of a Lipschitz activation makes these uniformly
dense on the compact input cube. The enlarged atom classes are only a
capacity bound; their finite-sample closures are compact subsets of the
finite-dimensional sample cube.

For the Gaussian comparison, set $E=d(d+1)B^2+2Hm^2B^2$ and $U=T^H(1+E)^{L/2}$. The perturbation proof in Lemma~\ref{lem:bv-comparison} uses only Lipschitzness and the activation offset, so it applies with this $U$, $B_0=1+\beta+2A(\sqrt E+1)$ and $K_g=2LB_0^H$. Its free-energy cost is
\[
 A_{g,n}=E+1+2(d+H+1)\log[4K_g(U+1)n].
\]
A uniform prior-risk bound is $1+2(4A_0^2)^H$, with $A_0=\max\{1,A,\beta\}$. Use these constants and the displayed capacity bound in the abstract theorem. The countable-cap population argument of Theorem~\ref{thm:network-main} remains valid because the statistical constant includes the envelope factor $T^H$, and $(U+1)^2$ is at most a fixed multiple of $T^{2H}(1+E)^L$.

The full-energy Gaussian posterior argument therefore applies with
$r=L$ to every globally Lipschitz activation, without a condition on
its derivative or its value at zero. The additional
$\log^{3H/2}(2n)$ only changes logarithmic factors in the calibration
and oracle. At deep scales $\gamma\ge2$, and at every scale
$\gamma\ge1$ for one hidden layer, the same teacher comparison and
free-energy root give
\[
 R(\widehat f_n)\le R(g)+\widetilde O_{H,A,\beta,d}
   \left(A_g R(g)^{1-1/L}n^{-1/L}+A_g^L/n\right).
\]
The prior remains one fixed dense independent Gaussian architecture.
The bounded-variation derivative proof remains useful because it avoids
these additional depth-dependent logarithmic powers. Tanh's separate
bounded-activation argument can retain its sharper depth exponent.
